\documentclass[12pt]{article}
\usepackage[utf8]{inputenc}
\usepackage[T1]{fontenc}
\usepackage[brazil]{babel}
\usepackage[a4paper,margin=2.5cm]{geometry}
\usepackage{enumitem}
\usepackage{amsmath,amssymb}
\setlist[itemize]{noitemsep,topsep=4pt}
\title{Sprint 3 - S3\_11\\Álgebra Linear III: Determinantes}
\author{Matheus Moreli Durão Silva\\Matrícula 12621CBS217\\Cibersegurança - Bacharelado - Noturno}
\date{2026}
\begin{document}
\maketitle

\section*{Entregável}
Determinantes, inversas e invertibilidade. Este documento registra o objetivo do card, a solução desenvolvida e a forma de conferência. Os arquivos de código e as evidências citadas permanecem na mesma pasta deste documento.

\section*{O que foi estudado e aplicado}
\begin{itemize}
\item Determinantes 2x2 e 3x3
\item Matriz singular
\item Inversa e solução única
\item Cenário invertível e cenário singular
\end{itemize}

\section*{Como conferir}
O avaliador pode abrir os arquivos desta pasta, executar os exemplos quando houver código e comparar a saída com a explicação apresentada. Nos cards ligados ao Inventário de Ativos, a solução foi relacionada ao Trabalho 2 sem modificar requisitos já atendidos do CRUD.

\section*{Inversa e determinante}
Para $B=\left[\begin{smallmatrix}2&1\\1&1\end{smallmatrix}\right]$, $\det B=1$ e $B^{-1}=\left[\begin{smallmatrix}1&-1\\-1&2\end{smallmatrix}\right]$. Para $C=\left[\begin{smallmatrix}1&2&0\\0&1&1\\1&0&1\end{smallmatrix}\right]$, Sarrus dá $1+2=3$ e cofatores na primeira linha dão $1-2(-1)=3$. Determinante diferente de zero implica transformação sem perda de dimensão e inversa existente.
\section*{Singularidade e domínio do programa}
Tome $A$ 3 por 3 com diagonal zero e cada entrada fora da diagonal igual a $0,5$. Cada peso é menor que 1, mas cada linha soma 1. Como $(I-A)(1,1,1)^T=0$, $I-A$ é singular. O programa proíbe soma de pesos por linha maior ou igual a 1: assim $I-A$ fica estritamente diagonal dominante e não singular. A checagem numérica adicional evita apresentar risco sem solução única.
\section*{Conclusão}
O artefato foi organizado para demonstrar tanto o conceito estudado quanto sua aplicação prática. As evidências de execução só são consideradas válidas quando produzidas pela ferramenta correspondente.
\end{document}
